\section{Method}
\label{sec:method}

The experiment is a substitution. VBVR-Pro-Bench gives a video model a first frame and a prompt and asks for the video in which the task gets solved. We hand the same two inputs to a coding agent in a sandbox, ask it to \emph{write a program} that renders that video, and score the result with the benchmark's own evaluator, unchanged. Everything downstream of the video file is identical for both families of system.

\subsection{The Benchmark}
\label{sec:bench}

\textbf{VBVR-Pro-Bench.} We use the video setting of VBVR-Pro-Bench \citep{vbvrpro2026}. The benchmark has 100 task families, each backed by a procedural generator and a hand-written, task-specific evaluator; every family contributes 5 instances, for 500 in total. The 50 \emph{In-Domain} families have training data in the VBVR-Pro SFT and RL corpora, so a video model trained on VBVR-Pro has seen the family, though not the instance; the 50 \emph{Out-of-Domain} families are held out from every training split. The split measures generalisation for the video models it was designed around. For the coding agents, none of which were trained on VBVR-Pro, both halves are unseen.

\textbf{Inputs and target.} An instance consists of \texttt{first\_frame.png}, a $1024\times1024$ image of the initial scene, and \texttt{prompt.txt}, a natural-language description of what the video should show. The reference video is 16\,fps and runs from 10 to 282 frames (0.6 to 17.6\,s) depending on the family. It ships with its final frame and a \texttt{metadata.json} holding the symbolic ground truth; the evaluator reads these, and no system under test sees them.

\textbf{Categories.} The kit assigns every task family to one of five faculties: Abstraction (25 tasks), Perception (28), Spatiality (16), Transformation (12) and Knowledge (19); the label is carried in the evaluator's per-instance output and is the one we use. Category means are plain means over the instances of the tasks carrying the label.

\subsection{Coding Agents as Visual Reasoners}
\label{sec:protocol}

\textbf{What the agent sees.} Each instance runs in a fresh Docker container whose only mounted directory, \texttt{/app}, holds three files: \texttt{first\_frame.png}, \texttt{prompt.txt} and an \texttt{instruction.md} generated by the harness. The instruction is the same for every agent and every model. Its essence is: \emph{``\texttt{first\_frame.png} is the first frame of a short video and \texttt{prompt.txt} describes what happens in it. Write code that produces the video. Deliver \texttt{output/video.mp4} (H.264, $1024\times1024$, 16\,fps, about $N$ frames) whose first frame is \texttt{first\_frame.png} and whose last frame shows the completed result, and keep the program at \texttt{solve.py} so it regenerates the video. The video is drawn by your code; do not call an image or video generation model. Change only what the prompt asks for; every other pixel stays as in the first frame. Pace the action over the full duration.''} The prompt is quoted verbatim inside the instruction. $N$ and the duration are probed from the reference video, so the agent knows the container spec the evaluator expects, as the video models do. Nothing else is disclosed.

\textbf{Sandbox.} The container image is \texttt{python:3.11-slim} with \texttt{ffmpeg}, NumPy, Pillow, OpenCV, imageio, SciPy and Matplotlib preinstalled, plus DejaVu, Liberation and Noto CJK fonts. The agent runs as an unprivileged user with 2 CPUs, 3\,GB of memory and no swap, a limit of 512 processes, and a wall-clock budget of 1{,}800\,s per instance, after which the harness kills the container and records no video. Network access is left on so that the agent CLI can reach its model API and \texttt{pip} can fetch a package; the rules forbid calling any generative image or video model, and a search of the closed agents' 1{,}500 event logs and of the 7{,}230 \texttt{solve.py} programs retained from the open-weight lanes finds no generation endpoint and no deep-learning framework; the packages agents did install were OCR and scientific-computing libraries. Each instance gets exactly one attempt: a container that exits without \texttt{output/video.mp4} is a failed instance, with no retry, no best-of-$n$ and no human intervention. All runs took place on one \texttt{c7i.4xlarge} CPU host between 12 and 16 September 2026.

\textbf{Agents.} An \emph{agent} here is an off-the-shelf command-line coding tool driving a model: it reads files, runs commands, edits code and iterates until it decides it is done. We invoke each tool through its official non-interactive entry point with approvals bypassed, and add no system prompt, skill or modification. Table~\ref{tab:agents} lists the three closed-model agents. Codex CLI receives the first frame as an image attachment as well as a file; the other agents open it from disk with their own file tools or inspect it programmatically, as they see fit. Claude Code is routed through Amazon Bedrock with the model identifier passed straight through. Codex CLI runs with its default reasoning effort; no reasoning-effort override was set for any lane.

\begin{table}[t]
\centering
\small
\setlength{\tabcolsep}{5pt}
\caption{The three closed-model coding agents. Each is the vendor's own command-line tool at the pinned version, run through its non-interactive mode with approvals bypassed. The 24 open-weight models are all driven by OpenCode 1.18.30 and are listed in Appendix~\ref{app:roster}.}
\label{tab:agents}
\begin{tabular}{lll}
\toprule
\textbf{Agent (CLI)} & \textbf{Version} & \textbf{Model} \\
\midrule
Codex CLI \citep{codexcli} & 0.154.0 & gpt-6-astra \\
Claude Code \citep{claudecode} & 2.1.268 & Claude Fable 5.1 (Amazon Bedrock) \\
Gemini CLI \citep{geminicli} & 0.59.0 & gemini-3.1-pro-preview \\
\bottomrule
\end{tabular}
\end{table}

\textbf{Open-weight roster.} To ask whether the result depends on a frontier closed model, we run the same protocol with OpenCode \citep{opencode}, an open-source agent that drives any model behind an OpenAI-compatible endpoint, over 24 open-weight models served by Amazon Bedrock: the GLM, Kimi, MiniMax, DeepSeek, Qwen, gpt-oss, Nemotron, Mistral, Llama and Gemma families (Appendix~\ref{app:roster}). The three Llama models are not served by Bedrock's OpenAI-compatible endpoint, and its streaming Converse API rejects tool use for them, so a small shim on the host accepts OpenCode's streaming request, calls the non-streaming Converse API and replays the answer as a short event stream; tool calls map one-to-one. DeepSeek R1 has no tool use on Bedrock in either mode, cannot act as a coding agent, and is listed as unsupported rather than scored. Every open-weight model receives the identical instruction, sandbox and budget. Several are text-only; the protocol does not require vision, since a program can read the pixels.

\subsection{Scoring}
\label{sec:scoring}

Every video is scored by the official VBVR-Pro-Bench kit, \texttt{run\_evaluation\_video.py}, without modification. The kit hands each instance to the family's evaluator, which decodes the video, compares it with the symbolic ground truth using classical computer vision (colour and shape primitives, Hungarian matching, multi-object tracking, OCR where the task shows text) and returns a score in $[0,1]$. We run it in \texttt{task\_specific\_only} mode, the mode behind the published leaderboard, on CPU. The score of an instance for which the agent produced no video is 0, exactly as for a video model that fails to generate. A lane's score is the mean over all 500 instances; In-Domain and Out-of-Domain means are over their 250 instances; category means follow the tag weighting of Section~\ref{sec:bench}. The video-model rows we compare against are the leaderboard numbers published with VBVR-Pro \citep{vbvrpro2026}, computed by the same kit; we do not rerun those models. One requirement matters: the kit must run under its pinned dependencies (NumPy~$<2$ plus its tracking and OCR packages), because under NumPy~2 four evaluators silently score every video, ground truth included, at 0; all numbers here come from a conforming environment (Appendix~\ref{app:evaluator-env}).

\subsection{Efficiency Accounting}
\label{sec:efficiency}

Each agent CLI emits a machine-readable event stream, which the harness keeps per instance as \texttt{events.jsonl} next to the container's exit code and wall time. From it we compute three quantities. \emph{Tool calls} count the actions the agent took: for Codex CLI, completed command-execution, file-change, MCP and web-search items; for Claude Code, \texttt{tool\_use} blocks in assistant messages; for Gemini CLI and OpenCode, their \texttt{tool\_use} events. \emph{Tokens} are the agent's own usage report summed over the run: Codex's per-turn usage; Claude Code's final result usage, with cache-creation and cache-read input counted as input; Gemini CLI's result statistics; OpenCode's per-step token record, with cached input counted as input and reasoning tokens counted as output. \emph{Wall time} is measured by the harness from container start to exit, so it includes model latency; we report the median over the 500 attempts. We also flag whether the agent ever called a tool, whether it hit the time limit, and whether a failed instance failed with or without tool use; these flags drive the failure anatomy in Section~\ref{sec:analysis}. Token counts are comparable within an agent and only indicative across agents, since tokenizers and caching differ.
